\documentclass[a4paper,11pt]{article}

% Packages
\usepackage[utf8]{inputenc}
\usepackage[T1]{fontenc}
\usepackage[french]{babel}
\usepackage{amsmath, amssymb, amsthm}
\usepackage{geometry}
\usepackage{xcolor}
\usepackage{tcolorbox}
\usepackage{enumitem}
\usepackage{fancyhdr}

% Configuration de la page
\geometry{margin=2cm}
\pagestyle{fancy}
\fancyhf{}
\fancyhead[L]{Bases de mathématiques}
\fancyhead[R]{État Vocabulaire++}
\fancyfoot[C]{\thepage}

% Environnements personnalisés
\tcbuselibrary{theorems}

\newtcolorbox{defbox}[1]{
    colback=blue!5!white,
    colframe=blue!75!black,
    fonttitle=\bfseries,
    title=#1
}

\newtcolorbox{propbox}[1]{
    colback=green!5!white,
    colframe=green!75!black,
    fonttitle=\bfseries,
    title=#1
}

\newtcolorbox{theorembox}[1]{
    colback=red!5!white,
    colframe=red!75!black,
    fonttitle=\bfseries,
    title=#1
}

\newtcolorbox{methodbox}[1]{
    colback=orange!5!white,
    colframe=orange!75!black,
    fonttitle=\bfseries,
    title=#1
}

\newtcolorbox{remarkbox}[1]{
    colback=purple!5!white,
    colframe=purple!75!black,
    fonttitle=\bfseries,
    title=#1
}

% Titre
\title{\textbf{Bases de mathématiques} \\ État Vocabulaire++}
\author{}
\date{}

\begin{document}

\maketitle
\tableofcontents
\newpage

\section{Assertions équivalentes}

\begin{defbox}{Définition : Assertions équivalentes}
Les assertions sont dites \textbf{équivalentes} lorsqu'elles possèdent la même valeur de vérité (même table de vérité), ainsi on note $P \Leftrightarrow Q$.
\end{defbox}

\subsection{Généralités sur les assertions mathématiques}

\begin{propbox}{Propriétés importantes}
\begin{itemize}
    \item Les assertions négatives sont notées $\overline{P}$, $\overline{P}$ ou encore $\neg P$
    \item On note $P$ ou $Q$ aussi par $P \vee Q$ (ou alors $P + Q$ en informatique)
    \item On note $P$ et $Q$ aussi par $P \wedge Q$ (ou alors $P \cdot Q$ en informatique)
    \item On note $P$ implique $Q$ ($P \Rightarrow Q$) si $P$ et $Q$ sont vraies, fausses ou si seulement $Q$ est vraie (car l'assertion dépend de $Q$)
    \item $(P$ et $Q) \Leftrightarrow (Q$ et $P)$ donc $(P$ ou $Q) \Leftrightarrow (P$ ou $Q)$
\end{itemize}
\end{propbox}

\subsubsection{Règles de logique}

\begin{propbox}{Associativité de "ou" et de "et"}
\begin{itemize}
    \item $(P$ ou $(Q$ ou $R)) \Leftrightarrow ((P$ ou $Q)$ ou $P) (P$ ou $(Q$ ou $R)) \Leftrightarrow ((P$ ou $Q)$ ou $R)$
    \item $\overline{(P \text{ ou } Q \text{ ou } R)} \Leftrightarrow (\overline{P} \text{ et } \overline{Q} \text{ et } \overline{R})$
\end{itemize}
\end{propbox}

\begin{propbox}{Distributivité de "ou" et de "et"}
\begin{itemize}
    \item $(P$ ou $(Q$ ou $R)) \Leftrightarrow ((P$ ou $Q)$ ou $(P$ ou $R))$
    \item $\overline{(P \text{ ou } (Q \text{ ou } R))} \Leftrightarrow ((\overline{P} \text{ et } \overline{Q}) \text{ et } (\overline{P} \text{ et } \overline{R}))$
\end{itemize}
\end{propbox}

\begin{propbox}{Transitivité de "et" sur le "ou"}
\begin{itemize}
    \item $(\text{non } (P$ ou $Q)) \Leftrightarrow (\text{non } P \text{ ou non } Q)$
    \item $(\text{non } (P$ et $Q)) \Leftrightarrow (\text{non } P \text{ ou non } Q)$
    \item $(P \Rightarrow Q \text{ et } Q \Rightarrow R) \Rightarrow P \Rightarrow R$
\end{itemize}
\end{propbox}

\subsection{Les raisonnements mathématiques}

\begin{methodbox}{Le raisonnement déductif (Syllogisme)}
On utilise la déduction : on dit que si $P$ et $P \Rightarrow Q$ sont vraies, alors $Q$ est vraie aussi.
\end{methodbox}

\begin{methodbox}{Le raisonnement par disjonction des cas}
On utilise par exemple si on peut montrer que $(P_1 \text{ ou } P_2) \Rightarrow Q$
\end{methodbox}

\begin{methodbox}{Le raisonnement par contraposée}
Pour rappel $(P \Rightarrow Q) \Leftrightarrow (\overline{Q} \Rightarrow \overline{P})$

Donc si l'on veut montrer que $P \Rightarrow Q$ est vraie, donc on montre que $\overline{Q} \Rightarrow \overline{P}$ est vraie. Pour démontrer qu'une assertion $P \Rightarrow Q$ est fausse, on fait l'hypothèse inverse, on suppose ou fausse en transposé un contre-exemple.
\end{methodbox}

\begin{methodbox}{Le raisonnement par l'absurde}
Pour démontrer que $\overline{Q} \Rightarrow \overline{P}$, pour montrer que $P \Rightarrow Q$ est vraie, on suppose l'hypothèse inverse, on fait une hypothèse absurde où on arrive à quelque chose qui est contradictoire.
\end{methodbox}

\begin{methodbox}{Le raisonnement par récurrence}
Il consiste à démontrer une assertion et fausse en transposé un contre-exemple qui soit contradictoire.

\textbf{Étapes :}
\begin{enumerate}
    \item[(A)] On définit $P(n)$
    \item[(2)] Initialisation : $P(n_0)$ est vraie
    \item[(3)] Hérédité : Soit $n \geq n_0$, on suppose $P(n)$ vraie et on montre $P(n+1)$
    \item[(4)] Conclusion : "Par le principe de récurrence, $P(n)$ est vraie pour tous $n$ entiers tels que $n \geq n_0$"
\end{enumerate}
\end{methodbox}

\newpage
\section{Le langage ensembliste et ses applications}

\subsection{Définitions}

\begin{defbox}{Définition : Ensemble}
Il est important d'écrire $E \in \mathbb{E}$ (interdit).

\textbf{Façon de définir un ensemble :}
\begin{itemize}
    \item Par extension (lister les éléments) : $E = \{1, 2, 5\}$
    \item Par union/compréhension (à partir d'une assertion) : $E = \{n - 1 | P(n \geq 1)\}$
\end{itemize}
\end{defbox}

\begin{defbox}{Définition : Élément d'un ensemble}
On note $E \subset F$ lorsque un ensemble $E$ est contenu dans un ensemble $F$ ($E$ est contenu dans $F$, $F$ contient $E$).

$E$ est une \textbf{partie} ou sous-ensemble de $F$.

$\forall x \in E$, on a $x \in E$ (\textbf{généralité}).
\end{defbox}

\begin{defbox}{Définition : Ensemble vide}
Si $E \subset F$ et $F \subset E$ alors $E = F$ (\textbf{anti-symétrie}).

Si $E \subset F$ et $F \subset G$ alors $E \subset G$ (\textbf{transitivité}) : $(E \subset F \text{ et } F \subset G) \Rightarrow E \subset G$.
\end{defbox}

\subsection{Le raisonnement par double inclusion}

\begin{methodbox}{Méthode}
Plutôt à 2 ensembles, pour montrer $E \subset G$, on montre $E \subset G$ et $G \subset E$.

Démontré il suffit $E \subset F$ puis, d'abord :
\begin{itemize}
    \item On note $A \in E$ ou $A \in G$ ($A$ avec $G$, $((E \cup F) \cup G = E \cup (F \cup G) = E \cup F \cup G)$
    \item On note $\overline{A} = E - A$ objets et objets de l'ensemble pour un couple d'objets. On a plusieurs $\{x, y\} = \{y, x\}$ ou seulement $\{x\} \in P$.
\end{itemize}

$\overline{A}$ signifie le complément de $A$ dans $E$, on a donc la collection des sous-ensembles des parties de $E$, $P(E)$, un ensemble des parties de $E$ : $P(E) \in \{\emptyset\}$ et $E \in P(E)$.

Soit $A \in P(E)$, on note \\
$\overline{E}$ : $\overline{A}$ de complémentaire de $A$ dans $E$
\end{methodbox}

\subsection{Opérations sur les ensembles}

\begin{defbox}{Définition : Intersection et union}
Soit $A$ et $B$ deux ensembles, on note :

\textbf{L'intersection :} $A \cap B = \{x : x \in A \text{ et } x \in B\}$

On note $\text{Card}(P(E))$ où $P$ et $Q$ : fini ex card $(P(E)) = \text{Card}(E)$.

Soit $P(E)$, on a donc la transitivité : $A \in P(E)$ et fini ex card $(P(E)) = \text{Card}(E)$.
\end{defbox}

\subsection{Relation binaire, Représentation schématique}

\begin{defbox}{Définition : Relation binaire}
Une relation binaire est une relation entre 2 ensembles $E$ et $E'$ telle que :

$E \times E' = \{(e, e') : e \in E, e' \in E'\}$ respecte donc la propriété :

\begin{itemize}
    \item \textbf{Réflexivité :} $\forall a \in E, \; a \mathcal{R} a$
    \item \textbf{Symétrique :} $\forall (a, b) \in E^2, \; (a \mathcal{R} b \Rightarrow b \mathcal{R} a) \Leftrightarrow (a \mathcal{R} b)$
    \item \textbf{Anti-symétrique :} $\forall (a, b) \in E^2, \; (a \mathcal{R} b \text{ et } b \mathcal{R} a) \Rightarrow a = b$
    \item \textbf{Transitive :} $\forall (a, b, c) \in E^3, \; (a \mathcal{R} b \text{ et } b \mathcal{R} c) \Rightarrow a \mathcal{R} c$
\end{itemize}

\[\Gamma = \{(a, b) \in E \times E : a \mathcal{R} b\} \subset E \times E, \; a \mathcal{R} b = \{(a, b) \in E \times E\}\]
\end{defbox}

\subsection{Relation d'équivalence, classe d'équivalence}

\begin{defbox}{Définition : Relation d'équivalence}
Une relation binaire dans $\mathbb{R}$ dans $E$ ($\mathbb{R} \Rightarrow E$) est appelée relation d'équivalence si $\mathcal{R}$ est réflexive, symétrique et transitive.

Soit $\mathcal{R}$ une relation d'équivalence dans $E$. On appelle classe d'équivalence de $a$, l'ensemble noté $\text{cl}(a)$ des éléments de $E$ en relation avec $a$.

\[\overline{a} = \{x \in E : a \mathcal{R} x = E\}\]

$\mathcal{VA} \mathcal{B}$ de $P(E)$, $A \mathcal{R} B \Leftrightarrow A \subset \mathbb{R}$ (Part avec dense d'équivalence)

$\mathcal{VA} E \subset P(E)$, $A \mathcal{R} B \Leftrightarrow A \subset E$ ($\mathbb{R}$ n'est pas une relation)

Si $E$ est un ensemble et $\mathcal{R}$ de relation d'équivalence dans $E$, alors pour $A \subset P(E)$.

\[A \mathcal{R} B \Leftrightarrow A = B \text{ et } \text{cl}(A) = \{A\}\]
\end{defbox}

\newpage
\section{Les relations d'ordres}

\begin{defbox}{Définition : Relation d'ordre}
Une relation binaire dans $\mathbb{R}$ dans $E$ ($\mathbb{R} \Rightarrow E$) est appelée relation d'ordre si $\mathcal{R}$ est réflexive, anti-symétrique et transitive.
\end{defbox}

\begin{defbox}{Définition : Ordre}
On coiffe $E$, $\mathcal{R}$ est un ensemble et $\mathcal{R}$, une relation d'ordre est appelée un \textbf{ensemble ordonné} :

\begin{itemize}
    \item $\forall A, B \in P(E)$, $A \mathcal{R} B \Leftrightarrow A \subset B$ est une relation d'ordre
    \item $\forall A, E \subset P(E)$, $A \mathcal{R} B \Leftrightarrow A \subset E$ est une relation d'ordre
    \item Par exemple, dans $\mathbb{Z}$, la relation "$\leq$" définie par :
    $\ell = \{(m, n) \in \mathbb{Z}^2 : m \leq n \Leftrightarrow m \in \mathbb{N}\}$
\end{itemize}
\end{defbox}

\subsection{Propriétés des relations d'ordre}

\begin{propbox}{Propriétés}
\textbf{$\mathcal{R}$ est-elle réflexive ?} $\forall m \in \mathbb{Z}$ alors $m - m = 0$ donc $m - m \in \mathbb{N}$

Soit $m \in \mathbb{N}$, alors $m - m = 0$ donc $m \leq m$

\textbf{$\mathcal{R}$ est-elle symétrique ?} $\forall (m, n) \in \mathbb{Z}^2, (m \mathcal{R} n) \in \mathbb{Z}^2, (m \mathcal{R} n) \cap P \mathcal{R} m$

Soit $m \in \mathbb{R}$, $\mathbb{R}$ sous $n - m \in \mathbb{N}$ et si on suppose que $m \mathcal{R} n$ puis

On dans $m - n \in \mathbb{N}$ soit que $k - P \in \mathbb{N}$

On peut donc observer que $m \mathcal{R} k$ et $k - m = (k - n) + (n - m) \in \mathbb{N}$

Donc $\mathcal{R}$ est transitive dans $\mathcal{R}$ est une relation d'ordre.
\end{propbox}

\newpage
\section{Applications, images, antécédents}

\subsection{Généralités}

\begin{defbox}{Définition : Application}
Soit 2 ensembles $E$ et $F$, on une relation $\mathcal{R} \subset E \times F$.

On peut définir une application $\sim$ et déterminer le domaine :
\begin{itemize}
    \item $\forall a \in E, \; \exists g \in F | a \mathcal{R} g$ : "chaque élément de $E$ doit avoir au moins une image dans $F$"
    \item $\forall a \in E, (a, g) \in \mathcal{R}$ et $a \mathcal{R} g$ où $a \mathcal{R} g'$, donc $g = g'$ : "chaque élément ne peut avoir qu'une seule image dans $F$"
\end{itemize}

Donc $\Gamma = \{(x, f(x)) \in E \times F | \forall x \in F, \; \{f(x)\} \in \text{Card}(g, g)\}$
\end{defbox}

\subsection{Image directe, Image réciproque}

\begin{defbox}{Définition}
Soit $E$ et $F$ 2 ensembles et $f : E \to F$. Soit $A$ une partie de $E$.

\textbf{L'image directe} de $A$ et le sous-ensemble de $F$ noté $f(A)$ :
\[f(A) = \{g \in F | \exists a \in A, \; g = f(a)\} | a \in A | f\]

\textbf{Remarquons :} On sait que $f(E) \subset F$ (à doigt image directe)

En montre aussi que $f$ et $g \in F$ par $g \in E$, $y = f(A)$ par suit directe de $F = y = f(E)$ et par double inclusion $f(E) = F$.

\textbf{L'image réciproque} de $A$ (ou la préimage de $A$) est le sous-ensemble de $E$ noté $f^{-1}(A)$ : $f^{-1}(A) = \{a \in E | f(a) \in A\}$

"L'ensemble des antécédents de $F$"
\end{defbox}

\begin{remarkbox}{Remarque}
Soit $E$, $F$, $f : E \to F$ Soit $\#$ en $F$.

Il ne faut pas confondre :
\begin{itemize}
    \item $f^{-1}(\{y\})$ qui n'est et le sous-ensemble $\{x\}$ de $F$ ou tout élément $y$ son un sur l'égalité du sous-ensemble de $E$
    \item $f^{-1}(y)$ qui n'a de sens que si $f$ est bijective, donc $f$ inverse à $f$ et $f^{-1}(y) \in E$ dans
\end{itemize}
\end{remarkbox}

\subsection{Dans le cas de $F$ est bijective}

\begin{propbox}{Propriété}
Si $y \in G : F$, alors $f^{-1}(y) = f^{-1}(y)$ c'est unique intérieure $y$ par $f$.
\end{propbox}

\begin{propbox}{Propriété : $\emptyset$ A est une partie de $F$}
Si $\emptyset$ $A$ est une partie de $F$ :
\[f^{-1}(\{y\}) = f^{-1}(y)\]

\textbf{Réciproquement :}
$\forall y \in F$

Si $A$ est une partie de $F$ : $f^{-1}(A) = \{f^{-1}(y)\}$
\end{propbox}

\newpage
\section{Arithmétique dans $\mathbb{Z}$}

\subsection{Propriétés de base}

\begin{propbox}{Propriétés}
Posons $a, b \in \mathbb{Z}$, $k \in \mathbb{Z}$, $a = k \cdot b$ où $b | a \rightarrow a = k \cdot b$

$|b| a \neq b/a \neq 1| \neq 1 | (2k + 1) \neq |a| (a | \text{cas})$
\end{propbox}

\begin{propbox}{Propriétés}
Quelques propriétés : Soit $a \in \mathbb{Z}$, $1/a$ est unique si $a$ et 0 (réfléchissez) et $0$ est possible

\[
\begin{cases}
Si \; 0 | a \text{ alors } a = 0 \\
Si \; a | 1 \text{ alors } a = \pm 1 \\
\end{cases}
\]

Soit $a, b, c \in \mathbb{Z}$ : $\begin{cases} 
a | a \text{ et } a | b \text{ alors } \forall (\alpha, \beta) \in \mathbb{Z}^2, \; a | (\alpha a + \beta b) \text{ (transitivité)} \\
a | b \text{ et } a | c \Rightarrow a + \beta c \\
a | b \Rightarrow (a | bc \text{ et } 0 \leq a) \Rightarrow (a | b \text{ mais } a \neq b) \\
(a | c \text{ et } 0 \leq a) \Rightarrow (a | b \text{ et } 0 | a) \\
(a | b \text{ et } b | a) \Rightarrow |a| = |b| \text{ (antisymétrie)} \\
\end{cases}$
\end{propbox}

\subsection{PGCD de 2 entiers}

\begin{defbox}{Définition : PGCD}
$M(a) = \{ka | k \in \mathbb{Z}\}$

$M(3) = \{3k | k \in \mathbb{Z}\}$ et $M(5) = \{5k\}$

Donc $\text{Atem}(a, b) = \min(M(a) \cap M(b) \cap \mathbb{N}^*)$.

$M(3) = 0$ et $M(f_n) = 1 \rightarrow \text{pgcd}(2, 3) = 6$

\textbf{On montre une propriété que nous ne vais à savoir à appli} :
\end{defbox}

\subsection{Les nombres premiers}

\begin{defbox}{Définition : Nombre premier}
Critères premiers entre eux : $(a, b) \in \mathbb{Z} : a$ et $b$ sont "premiers entre eux"

\begin{itemize}
    \item Si $a = 0$ et $b = ab$ avec $\text{pgcd}(a, b) = 1$, $\text{pgcd}(a, b) = 1$
    \item Si $a^2$ et $b^2$ sont "premiers entre eux"
\end{itemize}
\end{defbox}

\subsection{PPCM de 2 entiers}

\begin{defbox}{Définition : PPCM}
$M(a) = \{ka | k \in \mathbb{Z}\}$

$M(3) = \{3k | k \in \mathbb{Z}\}$ (donc $M(1) \subset \mathbb{N}(5) \subset M(\text{ppcm}(a, b))$)

Donc $\text{ppcm}(a, b) = \min(M(1) \cap M(5)) = \emptyset$.

On montre une propriété que nous ne vais à red.
\end{defbox}

\newpage
\section{La détermination en facteurs premiers}

\subsection{Comme l'Euclide}

\begin{methodbox}{Méthode}
Par $P$ premier et $> a \in \mathbb{Z}$

$P | a$ si $a$ dans $P | a$ ou $P | b$
\end{methodbox}

\begin{methodbox}{Démonstration par disjonction des cas}
\textbf{1er cas :} $P | a$ alors $P | a$ si $a$ sauf

\textbf{2ème cas :} $P \nmid a$
\end{methodbox}

\vspace{1cm}
\begin{center}
\textit{--- Fin du résumé ---}
\end{center}

\end{document}
